\section{Method}
\label{sec:method}

\subsection{Compiler}
\label{sec:compiler}

The compiler turns one (state, question) pair into a token sequence. The prompt has a fixed
system message, identical across backbones and identical to the one used by the project's
GLM direct-logit readout, and one user message of compact JSON,
\begin{center}
\small
\texttt{\{"state":}$S$\texttt{,"question":\{"type":}$\tau$\texttt{,"instructions":}$I$\texttt{,}\\
\texttt{"candidates":[\{"identifier":}$\iota_1$\texttt{,"meaning":}$m_1$\texttt{\},\dots]\}\}}
\end{center}
where $S$ is the state as JSON, $\tau$ the question type, $I$ the instructions and $m_i$
the meaning of candidate $i$: for a Choice the name, or the name and description together;
for a Noul the descriptions of \emph{false} and \emph{true}; for a Score the level text.
Appendix~\ref{app:prompt} gives the exact shape. Only the chat template and the tokenizer
change between backbones. Templates are rendered with the model's own chat template in a
sandboxed Jinja environment with thinking disabled (\texttt{enable\_thinking=False}); for
DeepSeek-V4.1-Flash, which ships no template, we use its official encoding reference in
chat mode. The compiler checks that the rendered template keeps exactly one insertion point
for the state and one for the question and that it keeps the system text, and it verifies
the checksums of the pinned tokenizer and template files before use.

\subsection{Candidate identifiers}
\label{sec:identifiers}

Candidate $i$ is presented with identifier $\iota_i$, which must be one ordinary token of the
backbone's vocabulary (not an added or control token), must decode back to itself, and must
differ from every other identifier's token. The preference order is the list of 255
identifier strings verified for the GLM readout, skipping any that is not a single token,
followed by a fixed extension of Greek, Cyrillic, Hebrew, Latin-1, Armenian and Georgian
letters. The GLM list starts with \texttt{A}--\texttt{Z}, \texttt{a}--\texttt{z} and
\texttt{0}--\texttt{9} and continues with multi-digit numerals from \texttt{10} onward.
Tokenizers that split digits, including those of the Qwen and Gemma models we tested, share
only those first 62 identifiers with GLM. Questions with at most 62 candidates therefore use
exactly GLM's identifiers, while the 77-, 200- and 255-candidate title questions use
extension letters. DeepSeek-V4.1-Flash, A.X-4.0-Light, Phi-4-mini-instruct and Tri-7B encode
all 255 GLM strings as single tokens and used them throughout. Section~\ref{sec:leaderboard}
shows why this matters.

\subsection{Untrusted text and the control vocabulary}
\label{sec:datatok}

The compiler uses two tokenizers built from the same files. The host tokenizer encodes the
template segments and may produce control tokens. The data tokenizer is the same vocabulary
and merge table with the list of added tokens removed, so no substring of the state,
instructions or candidate text can be encoded as a control token, including added tokens
that are not marked special (such as a thinking delimiter). Every data encoding is checked
again for reserved token IDs. This boundary is a structural guarantee about token IDs; it is
not a defence against semantic prompt injection.

\subsection{Piecewise candidate encoding}
\label{sec:piecewise}

When the question JSON is encoded in one call, merges can cross candidate boundaries and the
last token of a candidate is not well defined. In piecewise mode the compiler encodes the
question head, each candidate object and each separator separately and asserts that the
concatenated text is byte-identical to the compact JSON payload. This yields exact
candidate end positions, which the pointer head in Section~\ref{sec:pointer} needs. All
training, development, final and serving runs of \expone{} and its baseline use piecewise
encoding. The backbone comparison in Section~\ref{sec:leaderboard} used whole-payload
encoding; on the same development split Qwen3.8-27B scores 85.8\% whole-payload and 85.5\%
piecewise. The text is identical in both modes; only the token segmentation differs.

\subsection{First-position readout}
\label{sec:readout}

Let $h\in\mathbb{R}^d$ be the final hidden state at the last input position, which is the
position at which the assistant's answer would begin, and let $W$ be the output embedding
matrix. For offered identifier tokens $t_1,\dots,t_K$ the model's output is
\begin{equation}
  z_i = W_{t_i}^{\top} h, \qquad p = \operatorname{softmax}(z/T),
\end{equation}
computed in float32 from BF16 weights, with calibration temperature $T$. Nothing is sampled.
For diagnosis we also record the log of the full-vocabulary probability mass that falls on
the offered identifiers. For a Noul, the first two identifiers denote \emph{false} and
\emph{true} in that order; for a Score, identifiers follow the level order. In an earlier check on the GLM
path with real weights, 3 of 288 branches had a full-vocabulary argmax outside the offered
candidates; because only candidate rows are read, that token was never sampled, decoded or
returned.

\subsection{Backbone}
\label{sec:backbone}

The backbone is \texttt{Qwen/Qwen3.8-27B} at revision \texttt{1d4bf0f2ff60}
(Apache-2.0) \citep{qwen38_27b}. It has 27.8B parameters in a hybrid stack whose
linear-attention layers (Gated DeltaNet, \citealp{yang2024gdn}) and full-attention layers
alternate in a 3:1 ratio. We keep its original tokenizer and chat template and load BF16
weights. It was chosen on the development split (Section~\ref{sec:leaderboard}).

\subsection{LoRA adapter}
\label{sec:lora}

We attach LoRA \citep{hu2021lora} with rank $r=16$, $\alpha=32$ and no dropout to every
linear projection of the full-attention layers (\texttt{q\_proj}, \texttt{k\_proj},
\texttt{v\_proj}, \texttt{o\_proj}), the Gated DeltaNet layers (\texttt{in\_proj\_qkv},
\texttt{in\_proj\_z}, \texttt{in\_proj\_a}, \texttt{in\_proj\_b}, \texttt{out\_proj}) and
the MLPs (\texttt{gate\_proj}, \texttt{up\_proj}, \texttt{down\_proj}). The LM head and
embeddings stay frozen, and any vision, merger or multi-token-prediction modules are
excluded. The adapter has 116,727,808 trainable parameters; the base weights are frozen.
The adapter changes only the hidden state $h$; the readout rows $W_{t_i}$ are the original
ones.

\subsection{Objectives}
\label{sec:objectives}

For one question with candidate logits $z$, $p=\operatorname{softmax}(z)$ and gold index
$y$, we compare three losses.
\begin{align}
  \mathcal{L}_{\text{CE}} &= -\log p_y, \\
  \mathcal{L}_{\text{Brier}} &= \tfrac12 \textstyle\sum_i \left(p_i - \mathbb{1}[i=y]\right)^2, \\
  \mathcal{L}_{\text{RL}} &= -\tfrac14 \textstyle\sum_{j=1}^{4} (r_j - b_j)\log p_{a_j},
  \quad a_j \overset{\text{iid}}{\sim} p,\;
  r_j = \mathbb{1}[a_j=y] - \operatorname{sg}(p_{a_j}),\;
  b_j = \tfrac13 \textstyle\sum_{k\neq j} r_k ,
\end{align}
where $\operatorname{sg}$ stops the gradient. The RL learner receives only whether each
sampled action was correct. The reward $c_a - \operatorname{sg}(p_a)$ with independent action
samples and a leave-one-out baseline follows Anthony Maio's \texttt{eve-rlcd}
\citep{maio2026evercld} (MIT); our retention objective, distributed execution and evaluation
are a different training setup, and no \texttt{eve-rlcd} checkpoint is used.

\begin{proposition}
The expected gradient of $\mathcal{L}_{\text{RL}}$ equals the gradient of
$\mathcal{L}_{\text{Brier}}$.
\end{proposition}
\begin{proof}
The baseline $b_j$ depends only on the other samples, which are independent of $a_j$, and
$\mathbb{E}_{a\sim p}[\nabla\log p_a] = \sum_a \nabla p_a = 0$; so the baseline term has
zero expectation. For the remaining term, with $c_a=\mathbb{1}[a=y]$,
$\mathbb{E}_{a\sim p}\!\left[(c_a - p_a)\nabla\log p_a\right] = \sum_a (c_a-p_a)\nabla p_a
= -\nabla \tfrac12\sum_a (p_a - c_a)^2$, because $c_a$ does not depend on the parameters.
Negating gives $\mathbb{E}[\nabla\mathcal{L}_{\text{RL}}] = \nabla\mathcal{L}_{\text{Brier}}$.
\end{proof}

The RL arm therefore tests a sampled, higher-variance estimator of the half-Brier objective,
not a new source of information. Rows whose supervision is an observed mean rating
$\bar{s}$ (for example, STS and HelpSteer means) use the normalized squared error of the
expected level under every objective,
$\left(\left(\sum_i i\,p_i - \bar{s}\right)/(K-1)\right)^2$. This term fits a mean, not a
distribution.

\subsection{Residual pointer head (ablation)}
\label{sec:pointer}

As an architectural alternative to the vocabulary readout we add a residual bilinear score
between the final hidden state and the hidden state at each candidate's last token $e_i$:
\begin{equation}
  z_i \leftarrow z_i + s\,\frac{(Q h)^{\top} (K_p\, h_{e_i})}{\sqrt{256}},
  \qquad Q, K_p \in \mathbb{R}^{256\times d},
\end{equation}
with a learned scalar $s$. The pointer could correct identifier and list-position effects
without depending on how identifiers tokenize. With $s=0$ at initialization the first step
equals the vocabulary readout; Section~\ref{sec:limits} reports a retest that starts from
$s=1$.

\subsection{Temperature calibration}
\label{sec:calibration}

One scalar temperature $T$, shared by the three primitives, is fitted on the calibration
split (1,605 decisions) by bounded minimization of the mean negative log-likelihood over
$\log T\in[-3,3]$, and then applied unchanged to development and final data. It does not
change any argmax. The fitted values are $T=1.1489$ for \expone{} and $T=1.2596$ for the
zero-shot baseline.

\subsection{Batching and shared-prefix branching}
\label{sec:branching}

The questions of one request are compiled into separate sequences and run as one
right-padded batch with explicit position IDs. Each sequence is read at its own last
position; because both causal attention and the causal linear-attention recurrence run left
to right, padding placed after that position cannot affect it. Batches are split into
chunks of at most 131,072 padded tokens.

In shared-prefix mode the longest token prefix common to all of a request's sequences (the
template, the state and the start of the question object) is prefilled once with a cache.
For each chunk of suffixes the cache is deep-copied and expanded along the batch dimension,
covering both the full-attention keys and values and the Gated DeltaNet convolution and
recurrent states. The suffixes then run from the copy with positions offset by the prefix
length; the parent cache is never modified. We checked, on a small random hybrid model with
the GPU kernels and on the real model, that the implementation we use (transformers 5.17
with flash-linear-attention kernels) carries the linear-attention state through a
multi-token continuation. The API server's transformers engine uses this mode when the common
prefix is at least 1,024 tokens. Section~\ref{sec:determinism} measures its numerical agreement with full
sequences.

\subsection{Usage accounting}
\label{sec:usage}

\texttt{usage.input\_tokens} counts the request's own content as the backbone tokenizer
encodes it: the state once, plus each question's instructions, candidates and JSON
punctuation. The fixed template and the per-question copies of the state are not counted.
The number of tokens actually processed is reported separately in a response header. This
counter is Bobcat's own definition and is not claimed to match TypeSafe's token counter.
